\documentclass[conference]{IEEEtran}

% ── Packages ──────────────────────────────────────────────────────────────
\usepackage[T1]{fontenc}
\usepackage[utf8]{inputenc}
\usepackage{graphicx}
\usepackage{amsmath}
\usepackage{amssymb}
\usepackage{booktabs}
\usepackage{array}
\usepackage{multirow}
\usepackage{url}
\usepackage{hyperref}
\usepackage{cite}
\usepackage{xcolor}
\usepackage{balance}
\usepackage{float}
\usepackage{caption}
\usepackage{subcaption}
\usepackage{microtype}

\hypersetup{
  colorlinks=true,
  linkcolor=blue,
  citecolor=blue,
  urlcolor=blue
}

% ── Title block ────────────────────────────────────────────────────────────
\title{Applying Reinforcement Learning to the\\Connection Game Y}

\author{
  \IEEEauthorblockN{Advait Shinde\IEEEauthorrefmark{1},
                    Jasper Morgal\IEEEauthorrefmark{2},
                    Nickzad Bayati\IEEEauthorrefmark{3},
                    Toney Zhen\IEEEauthorrefmark{4}}
  \IEEEauthorblockA{Department of Computer Engineering\\
    San José State University, San José, CA\\
    \{advait.shinde, jasper.morgal, nickzad.bayati, toney.zhen\}@sjsu.edu}
    }


\begin{document}
\maketitle

% ── Abstract ───────────────────────────────────────────────────────────────
\begin{abstract}
Y is a connection strategy game played on a triangular hexagonal board where the
goal is to form a connected group of stones touching all three sides. While
strategically complex games such as Go and Hex have been extensively explored with
reinforcement learning, Y remains largely unexamined in the literature. In this
work, we implement and evaluate nine distinct agents across three paradigms:
tree-search agents (MCTS and policy/value-guided MCTS variants), temporal
difference learning agents (TD(0), TD($\lambda$), and TD-CNN), and
Stable-Baselines3 deep RL agents (MaskablePPO, PPO, A2C, and DQN). All agents
compete in round-robin tournaments evaluated via Bradley--Terry Elo ratings.
Our results show a clear three-tier hierarchy: search-based agents dominate
($\text{Elo} > 1869$), SB3 agents with action masking are competitive with
hand-crafted heuristics ($\text{Elo} \approx 1740$), and pure value-based
TD agents form a weaker tier ($\text{Elo} < 1250$). Among deep RL methods,
action masking proves to be the single most impactful design choice, enabling
MaskablePPO to outperform all TD agents and match the domain-knowledge
heuristic after only 2\,M training steps. We also demonstrate that iterative
self-play curricula benefit MCTS-guided methods but can destabilize on-policy
gradient methods such as PPO.
\end{abstract}

\begin{IEEEkeywords}
reinforcement learning, Monte Carlo tree search, temporal difference learning,
connection games, policy gradient, action masking, self-play, Elo rating
\end{IEEEkeywords}

% ══════════════════════════════════════════════════════════════════════════
\section{Introduction}
% ══════════════════════════════════════════════════════════════════════════

Reinforcement learning (RL) has demonstrated remarkable success in combinatorial
board games, from TD-Gammon \cite{tesauro1995} and Deep Blue
\cite{campbell2002} to AlphaGo \cite{silver2016} and AlphaZero
\cite{silver2017}. These advances have inspired the application of similar
methods to a wide range of abstract strategy games. Y, invented by John Nash and
Craige Schensted in the 1950s, is a two-player zero-sum connection game played on
a triangular hexagonal board. The first player to form a connected group of their
stones that simultaneously touches all three sides of the board wins. Y shares
structural properties with Hex: determinism, perfect information, no chance
elements, and no draws. However, its three-sided win condition gives rise to
qualitatively different strategic patterns that have not yet been studied in the
RL literature.

This paper makes the following contributions:

\begin{enumerate}
  \item A complete open-source implementation of Y as a Cython-accelerated game
    engine with a Gymnasium-compatible environment. (\href{https://github.com/Jasper-256/y-ai}{https://github.com/Jasper-256/y-ai})
  \item Training and evaluation of nine agents spanning three paradigms: MCTS
    variants, TD learning, and SB3 deep RL.
  \item A comprehensive round-robin tournament with Bradley--Terry Elo analysis
    identifying a clear performance hierarchy.
  \item Empirical evidence that action masking is the dominant factor separating
    effective from ineffective deep RL agents on this domain.
  \item An analysis of self-play curriculum effects, showing divergent outcomes
    for search-based versus gradient-based methods.
\end{enumerate}

% ══════════════════════════════════════════════════════════════════════════
\section{Related Work}
% ══════════════════════════════════════════════════════════════════════════

\subsection{Search-Based Methods}

AlphaZero \cite{silver2017} established the modern paradigm for game-playing AI
by combining MCTS with jointly trained policy and value networks in a self-play
loop. The policy network provides PUCT-guided priors that focus search on
promising moves, while the value network replaces Monte Carlo rollouts with
direct position evaluation. MuZero \cite{schrittwieser2021} extended this
framework to settings where game dynamics are not explicitly known by learning a
latent model of the environment. Our PV-MCTS and SP-PV-MCTS agents follow the
AlphaZero design philosophy, but use lightweight NumPy MLPs rather than deep
residual networks due to hardware constraints.

\subsection{Temporal Difference Learning}

TD-Gammon \cite{tesauro1995} demonstrated that TD($\lambda$) with a neural
network value function can reach expert-level play in Backgammon. TD(0) is a
special case that updates only from the immediately following state, while
TD($\lambda$) uses eligibility traces to propagate credit across an entire
episode. These methods have been applied to Chess \cite{veness2009} and other
board games, typically as value estimators for 1-ply or multi-ply lookahead.

\subsection{Deep RL for Board Games}

Anthony et al.\ \cite{anthony2017} proposed Thinking Fast and Slow, which
combines a fast neural network policy with a slower MCTS planner and showed
strong results on the game of Hex. Stable-Baselines3 \cite{raffin2021} provides
reliable, well-tested implementations of PPO \cite{schulman2017}, A2C, and DQN
\cite{mnih2015} that have been applied to a variety of sequential decision-making
problems. MaskablePPO from sb3-contrib \cite{sb3contrib} extends PPO with action
masking, preventing the policy from placing probability mass on illegal moves a
critical feature for board games with a shrinking legal action set.

\subsection{The Game of Y and Hex}

Hex has received considerable RL attention. HexHex \cite{harbecke2025} is an
AlphaGo Zero adaptation for Hex that demonstrates strong play through self-play
with policy/value networks. To our knowledge, no prior work has applied RL to
Y.

% ══════════════════════════════════════════════════════════════════════════
\section{Game Environment}
% ══════════════════════════════════════════════════════════════════════════

\subsection{Rules of Y}

Y is played on a triangular board of side length $N$, comprising
$N(N+1)/2$ hexagonal cells. Two players alternate placing stones of their color.
A player wins by forming a connected group of stones that simultaneously touches
all three sides of the triangle. Unlike Hex, all three sides are shared win
conditions rather than belonging to a specific player, which prevents draws and
creates richer tactical situations near board corners.

\subsection{Implementation}

The game engine is implemented in Cython for performance. The \texttt{Board}
class stores cell occupancy in a flat integer array and uses a BFS-based win
detector that checks whether a connected component of a given player's stones
touches all three sides. The \texttt{Game} wrapper tracks turn order, move
history, and terminal state.

All experiments use a board of side length $N = 7$, giving 28 cells and an
action space of $|\mathcal{A}| = 28$. Boards of size 9 was implemented on search based algorithms, but not for stable baseline algorithms, hence 9 to 19 board sizes are planned for
future evaluation but were not included in the current study due to
the rapid growth of computational requirements with board size.

\subsection{Gymnasium Environment}

We provide two Gymnasium-compatible environments. \texttt{YGameEnv} exposes a
\texttt{Discrete(28)} action space suitable for PPO, A2C, DQN, and MaskablePPO.
\texttt{ContinuousYGameEnv} exposes a \texttt{Box(-1,1, shape=(28,))} action
space for DDPG, TD3, and SAC, where the environment applies a legality mask to
the continuous output before taking an argmax to select the actual cell. The
observation is a float32 vector of length $28 \times 3 + 1 = 85$ encoding, for
each cell, whether it is empty, occupied by the current player, or occupied by the
opponent, plus a bias feature. This encoding is always computed from the
perspective of the player to move, ensuring that a single model can play as
either colour.

The environment calls the opponent's \texttt{choose\_move()} inside
\texttt{step()}, so the RL agent always receives states where it is its turn to
move. An illegal move immediately terminates the episode with reward $-1$; a win
returns $+1$ and a loss returns $-1$; all other transitions return $0$.

% ══════════════════════════════════════════════════════════════════════════
\section{Agent Architectures}
% ══════════════════════════════════════════════════════════════════════════

Fig.~\ref{fig:arena} provides an overview of the full system, showing the
relationship between agents, the Y-Hex environment, the SB3 training pipeline,
the arena round-robin engine, and the Elo aggregation step.

\subsection{Baseline Agents}

\textbf{Random Agent.} Selects a uniformly random legal move. Serves as the
floor baseline.

\textbf{Heuristic Agent.} Performs 1-ply lookahead and scores each successor
position by (i) how many distinct board sides the best-connected component of the
current player touches (primary), (ii) the stone count of that component, and
(iii) the number of empty cells adjacent to the player's stones. No learning
is involved; domain knowledge is encoded directly in the evaluation function.

\subsection{MCTS Agents}

\textbf{MCTS.} Standard UCB1-based Monte Carlo Tree Search implemented in Cython.
Each iteration performs selection, expansion, a random rollout to game end, and
backpropagation. All experiments use 500 iterations per move.

\textbf{PV-MCTS.} A policy/value-guided MCTS agent following the AlphaZero
architecture. A two-layer MLP (hidden size 192) with a shared trunk, policy head
(softmax over 28 cells), and value head (sigmoid scalar) is trained from 200
teacher self-play games generated by a strong MCTS agent (2\,000 iterations). At
inference, PUCT selection weights move priors from the policy network. The value
head is blended with terminal rollouts at leaf nodes (weight 0.25).

\textbf{SP-PV-MCTS.} The self-play variant of PV-MCTS. The same MLP is trained
entirely from its own MCTS search statistics over 7 generations of 60 games each
(300 search iterations per move), with no external teacher. Policy targets are the
visit-count distributions from MCTS; value targets are final game outcomes.

\textbf{SP-PV-CNN.} Replaces the MLP trunk of SP-PV-MCTS with a small CNN.
The triangular board is embedded diagonally into an $N \times N$ grid with four
input channels (exists mask, mine, opponent, empty), followed by three
$3 \times 3$ convolutional layers (24 channels each) and a small dense head.
Six-way board-symmetry augmentation is applied during training. The agent uses
350 MCTS iterations per move.

\textbf{SP-Policy-CNN.} A policy-only CNN that plays directly from the network
output with no MCTS at inference time. Trained by REINFORCE on its own self-play
trajectories with symmetry augmentation. Represents the lower bound of the
CNN-based approach family.

\subsection{Temporal Difference Agents}

All TD agents use 1-ply lookahead: evaluate every legal successor from the
opponent's perspective, then choose the move minimising the opponent's estimated
win probability.

\textbf{TD(0).} A two-layer MLP ($85 \to 128 \to 1$, ReLU/sigmoid) trained by
TD(0) self-play. The value function $V(s) \in [0,1]$ estimates the probability
that the current player wins. Training runs for 2\,000 self-play games with
$\epsilon$-greedy exploration ($\epsilon = 0.1$).

\textbf{TD($\lambda$).} Extends TD(0) with eligibility traces ($\lambda = 0.7$),
propagating TD errors across entire episodes. Uses the same network and training
budget.

\textbf{TD-CNN.} Replaces the MLP with the same diagonal-embedding CNN described
for SP-PV-CNN, but trained with TD(0) updates rather than MCTS self-play. Three
convolutional layers ($32, 32, 16$ channels) feed a 64-unit dense head.

\subsection{SB3 Deep RL Agents}

All SB3 agents use \texttt{MlpPolicy} and the default network architecture
provided by Stable-Baselines3. Training uses 500\,000 steps per stage with an
iterative self-play curriculum of three stages, for a total of 2\,M steps per
agent. In each self-play stage, the model from the previous stage is frozen via
save/load (not \texttt{deepcopy}, which is incompatible with PyTorch weight
normalisation) and used as the opponent.

\textbf{MaskablePPO.} PPO \cite{schulman2017} with action masking from
sb3-contrib. The policy's logits for illegal cells are set to $-\infty$ before
the softmax, ensuring zero probability mass on illegal moves at every gradient
step. This is the recommended algorithm for discrete board games.

\textbf{PPO.} Standard PPO without action masking. Illegal moves are penalised
with reward $-1$ and episode termination, requiring the policy to learn the
legality constraint from environment feedback alone.

\textbf{A2C.} Synchronous Advantage Actor-Critic with $n = 5$ environment steps
per update. Simpler than PPO (no clipping, no multiple epochs per rollout), which
trades peak performance for training stability.

\textbf{DQN.} Deep Q-Network \cite{mnih2015} with an experience replay buffer
(100\,000 transitions), $\epsilon$-greedy exploration (final $\epsilon = 0.05$),
and a target network. Off-policy learning provides better sample efficiency than
on-policy methods on small boards.



% ══════════════════════════════════════════════════════════════════════════
\section{Evaluation Methodology}
% ══════════════════════════════════════════════════════════════════════════

\subsection{Training Evaluation}

During TD agent training, checkpoint evaluations are run every 100 games against
a random baseline and a heuristic baseline. At each checkpoint, the agent plays
50 games as both P1 and P2 against each opponent (100 games per checkpoint per
opponent). Fig.~\ref{fig:training_curves} shows win-rate trajectories for a
representative TD(0) run (hidden size 256).

\subsection{Round-Robin Tournament}

All trained agents compete in a round-robin tournament. For each ordered pair
$(A_i, A_j)$ with $i \neq j$, $N_{\text{games}}$ games are played with $A_i$ as
Player 1 and $A_j$ as Player 2. Because Y has a first-player advantage (no swap
rule is implemented), both orderings are played and results are aggregated. The
primary evaluation uses $N_{\text{games}} = 50$ per ordered pair (100 per
unordered pair).

\subsection{Bradley--Terry Elo}

We fit a Bradley--Terry model \cite{bradleyterry1952} to the tournament results.
The strength parameter $\alpha_i$ for agent $i$ satisfies
\begin{equation}
  P(i \text{ beats } j) = \frac{e^{\alpha_i}}{e^{\alpha_i} + e^{\alpha_j}},
\end{equation}
which is equivalent to a logistic Elo model with scale factor
$\frac{400}{\ln 10} \approx 173.7$. Strengths are estimated by Newton--Raphson
on the log-likelihood. We fix the weakest agent's strength to zero, centre all
strengths, and report ratings on the standard Elo scale anchored at 1500.
Standard errors are derived from the inverse Fisher information matrix.

% ══════════════════════════════════════════════════════════════════════════

% ── Full-width architecture figure ─────────────────────────────────────────
\begin{figure*}[t]
  \centering
   \includegraphics[width=\textwidth]{arena.png}
  \caption{System architecture overview. \textbf{Left:} Agent roster covering
    SB3 deep RL agents (MaskablePPO, PPO, A2C, DQN) and MCTS and TD agents. \textbf{Centre:} The Y-Hex
    Gymnasium environment showing the triangular board, observation encoding,
    and action space with legality masking. \textbf{Right:} The SB3 training
    pipeline with optional iterative self-play stages and model checkpointing.
    \textbf{Bottom:} The Arena round-robin tournament engine (ordered-pair
    matchup table and single-game loop) feeding into Bradley--Terry Elo
    aggregation.}
  \label{fig:arena}
\end{figure*}

% ══════════════════════════════════════════════════════════════════════════
\section{Experiments and Results}
% ══════════════════════════════════════════════════════════════════════════

\subsection{Training Dynamics (TD Agents)}

Fig.~\ref{fig:training_output} shows the checkpoint win-rate trajectories for a
TD(0) agent (hidden size 256) trained over 2\,000 self-play games on a size-7
board, evaluated every 100 games against random and heuristic opponents.

Against the random baseline, the agent quickly rises to approximately 70\% win
rate within the first 200 games and continues improving to over 85\% by game
2\,000, demonstrating that meaningful learning occurs early in training. Against
the heuristic baseline, progress is slower and more volatile: the agent reaches
only 14--23\% win rate across all checkpoints, confirming the difficulty of
surpassing the hand-coded evaluation function with a shallow value-based approach.

The ``current vs.\ previous checkpoint'' panel reveals that the agent does not
monotonically improve over successive checkpoints: in 5 out of 20 checkpoints,
the current agent loses to the previous snapshot in one or both seat assignments.
This oscillation is characteristic of TD self-play and motivates more structured
training curricula.

\subsection{Overall Elo Leaderboard}

Table~\ref{tab:elo_full} presents the final Bradley--Terry Elo ratings from a
unified round-robin tournament combining all nine trained agents plus the
random and heuristic baselines. The SB3 agents were evaluated in a separate
tournament with identical methodology; their relative ordering is well-established
but their absolute Elo values are not directly comparable to the main tournament
due to different opponent pools.

\begin{table}[h]
  \centering
  \caption{Final Bradley--Terry Elo Leaderboard (all agents, $N=50$ games per
    ordered pair, board size 7). SB3 agents evaluated in a separate tournament;
    see Section~\ref{sec:sb3results} for details.}
  \label{tab:elo_full}
  \renewcommand{\arraystretch}{1.15}
  \begin{tabular}{clrrl}
    \toprule
    \textbf{Rank} & \textbf{Agent} & \textbf{Elo} & \textbf{$\pm$SE} & \textbf{Tier} \\
    \midrule
    1  & SP-PV-MCTS(500)   & 2085 & 42 & Search + Self-play \\
    2  & PV-MCTS(500)      & 2071 & 42 & Search + Self-play \\
    3  & SP-PV-CNN(500)    & 2018 & 41 & Search + Self-play \\
    4  & MCTS(500)         & 1869 & 38 & Pure Search        \\
    \midrule
    5  & MaskablePPO       & 1740 & 18 & SB3 RL             \\
    6  & Heuristic         & 1716 & 17 & Domain Knowledge   \\
    7  & DQN               & 1528 & 15 & SB3 RL             \\
    8  & A2C               & 1429 & 15 & SB3 RL             \\
    9  & PPO               & 1369 & 16 & SB3 RL             \\
    \midrule
    10 & SP-Policy-CNN     & 1227 & 31 & Policy CNN         \\
    11 & TD(0)             & 1224 & 31 & TD Learning        \\
    12 & TD($\lambda=0.7$) & 1211 & 31 & TD Learning        \\
    13 & TD-CNN            & 1065 & 33 & TD Learning        \\
    14 & Random            &  867 & 40 & Baseline           \\
    \bottomrule
  \end{tabular}
\end{table}

Three distinct performance tiers are clearly visible:
\begin{itemize}
  \item \textbf{Tier 1 (Elo $>$ 1869):} Search-based agents. The ${\sim}150$
    Elo gap separating MCTS(500) from the Heuristic baseline implies decisive advantage of runtime
    tree search over any learned policy alone.
  \item \textbf{Tier 2 (Elo 1300--1750):} SB3 RL agents and the Heuristic.
    MaskablePPO sits above the Heuristic; the remaining SB3 agents fall below
    it. All four SB3 methods substantially outperform the TD agents.
  \item \textbf{Tier 3 (Elo $<$ 1250):} TD agents, SP-Policy-CNN, and
    Random. The three TD variants cluster within 160 Elo points of each other
    and near the no-search CNN baseline, suggesting a shared representational
    ceiling for value-only approaches without lookahead.
\end{itemize}

\subsection{Search Agent Analysis}

\subsubsection{MCTS}
Plain MCTS with 500 iterations per move achieves Elo 1869, defeating every
non-search agent with near-100\% win rates. It loses to PV-MCTS and SP-PV-MCTS
with 30\% and 25\% win rates respectively, confirming that policy network priors
meaningfully improve search efficiency.

\subsubsection{PV-MCTS vs.\ SP-PV-MCTS}
PV-MCTS (trained from a teacher MCTS) achieves Elo 2071 and SP-PV-MCTS (pure
self-play) achieves Elo 2085 a difference of only 14 Elo points with
overlapping error bars ($\pm 42$). Head-to-head, SP-PV-MCTS wins 22/40 games,
suggesting a marginal edge for self-play training but no statistically significant
advantage over teacher distillation at this scale.

\subsubsection{SP-PV-CNN}
The CNN-trunked self-play agent achieves Elo 2018, placing third. Notably it
exhibits a strong first-player advantage: as P1 it wins 90\% of games against
SP-PV-MCTS, but loses 19/20 as P2. This asymmetry suggests the CNN learned an
aggressive opening strategy tuned to first-move advantage rather than a general
positional understanding.

\subsubsection{SP-Policy-CNN}
Without MCTS at inference time, the self-play CNN collapses to Elo 1227 nearly
800 Elo below SP-PV-CNN. The runtime search component contributes the vast
majority of the MCTS-family performance, not the network itself.

\subsection{TD Agent Analysis}

\subsubsection{TD(0) vs.\ TD($\lambda$)}
TD(0) (Elo 1224) and TD($\lambda = 0.7$) (Elo 1211) are separated by only 13
Elo points ($\ll$ the $\pm 31$ standard errors). Their head-to-head record is
perfectly one-sided (TD(0) wins all 40 games), but this reflects a degenerate
fixed-point in the deterministic policy rather than a genuine strength gap. The
two agents are functionally equivalent at the population level.

The minimal impact of eligibility traces corroborates the conclusion that the
bottleneck for TD agents in Y is the depth of lookahead, not the quality of
credit assignment. A single-step bootstrapping target is already sufficient to
extract most of the learning signal available to a 1-ply agent.

\subsubsection{TD-CNN}
The CNN value function achieves Elo 1065, underperforming both MLP-based TD
variants. This is counterintuitive given the richer representational capacity of
the CNN, but is consistent with prior observations that convolutional features
tuned to spatial patterns can overfit to specific board configurations during
self-play, reducing generalization.

\subsection{SB3 Deep RL Agent Analysis}
\label{sec:sb3results}

\subsubsection{Experimental Setup}
SB3 agents were trained on Google Colab (T4 GPU). Each agent was trained for 500\,000 steps per
stage across four stages (stage 0 vs.\ RandomAgent; stages 1--3 self-play),
totalling 2\,M steps per agent. The SB3 tournament used $N = 50$ games per
ordered pair with baselines Random and Heuristic; TD agents and MCTS are included
in the main tournament table.

Fig.~\ref{fig:sb3_elo} shows the Elo rating achieved by each SB3 algorithm as a
function of cumulative training iterations, illustrating both the benefit of
self-play stages and the divergence between masked and unmasked methods.

\subsubsection{MaskablePPO}
MaskablePPO achieves Elo 1740, outperforming the Heuristic (1716) and all TD
agents. The action mask eliminates gradient waste on illegal moves: with 28 cells
and an average of $\sim\!18$ legal moves per position, approximately one-third
of the standard PPO policy gradient flows toward cells that cannot legally be
played. Masking redirects this entirely toward learning a better policy over legal
moves. MaskablePPO wins 54/100 games against the Heuristic, a 54\% win rate
corresponding to the 24 Elo point gap.

\subsubsection{PPO}
Plain PPO achieves Elo 1369, 371 points below MaskablePPO and below both A2C and
the Heuristic. Notably, PPO was observed to perform better at 50\,000 steps
(Elo $\approx 1273$) than after the full 2\,M step curriculum, indicating policy
collapse during self-play. The on-policy clipping mechanism of PPO appears to
interact poorly with the non-stationary opponent distribution in self-play,
causing the policy to oscillate rather than converge. This effect is well
documented for on-policy methods in adversarial self-play settings \cite{balduzzi2019}.

\subsubsection{A2C}
A2C achieves Elo 1429, placing above PPO despite being the simpler algorithm.
Its synchronous update rule (no clipping, no multiple epochs per rollout) is more
conservative and less prone to the overshooting that destabilises PPO under
self-play. A2C wins 66/100 head-to-head against PPO, a striking reversal of
the typical relative ranking of these algorithms on single-task environments.

\subsubsection{DQN}
DQN achieves Elo 1528, the second-highest SB3 result. Its replay buffer provides
effective sample reuse that on-policy methods cannot exploit. However, DQN
exhibits a pronounced first-player advantage: it wins 50/50 as P1 against
MaskablePPO but 0/50 as P2, suggesting it learned a position-dependent opening
strategy rather than a symmetric game policy. This is a known failure mode for
Q-learning on asymmetric games without careful opponent randomisation.

\begin{table}[h]
  \centering
  \caption{SB3 Agent Elo ratings (separate tournament, $N=50$ per ordered pair).}
  \label{tab:sb3_elo}
  \renewcommand{\arraystretch}{1.15}
  \begin{tabular}{clrrc}
    \toprule
    \textbf{Rank} & \textbf{Agent} & \textbf{Elo} & \textbf{$\pm$SE} & \textbf{Win\%} \\
    \midrule
    1 & MaskablePPO & 1740 & 18 & 80.2\% \\
    2 & Heuristic   & 1716 & 17 & 77.6\% \\
    3 & DQN         & 1528 & 15 & 53.8\% \\
    4 & A2C         & 1429 & 15 & 40.4\% \\
    5 & PPO         & 1369 & 16 & 32.6\% \\
    6 & Random      & 1218 & 19 & 15.4\% \\
    \bottomrule
  \end{tabular}
\end{table}

\subsection{Head-to-Head Heatmap}

Fig.~\ref{fig:heatmap} shows the pairwise win rates (combined across both seat
assignments) for the main tournament. The clear block structure confirms the
three-tier hierarchy: each tier wins near-100\% against lower tiers and splits
roughly evenly within its tier.

\subsection{First-Player Advantage}

Table~\ref{tab:p1advantage} summarises first-player win rates for selected
matchups. Without a swap rule, P1 wins 60--100\% of games against similarly
skilled opponents. The strongest agents partially offset this by playing
differently as P2: SP-PV-MCTS as P2 wins 30\% against MCTS (P1), significantly
above random, indicating genuine second-player strategy. DQN's 0\% win rate as
P2 against MaskablePPO stands out as an extreme case of P1-dependent play.

\begin{table}[h]
  \centering
  \caption{First-player win rates for selected matchups ($N=50$ per direction).}
  \label{tab:p1advantage}
  \renewcommand{\arraystretch}{1.15}
  \begin{tabular}{llc}
    \toprule
    \textbf{P1 Agent} & \textbf{P2 Agent} & \textbf{P1 Win\%} \\
    \midrule
    SP-PV-MCTS & SP-PV-CNN     & 95\%  \\
    PV-MCTS    & SP-PV-MCTS    & 60\%  \\
    MCTS       & PV-MCTS       & 30\%  \\
    MaskablePPO & Heuristic    & 56\%  \\
    DQN        & MaskablePPO   & 100\% \\
    DQN        & MaskablePPO (P2) & 0\%  \\
    Heuristic  & Random        & 95\%  \\
    TD(0)      & Heuristic     & 60\%  \\
    \bottomrule
  \end{tabular}
\end{table}

% ── Figures referenced in results ─────────────────────────────────────────

\begin{figure}[h]
  \centering
  \includegraphics[width=\columnwidth]{training_output.png}
  \caption{Mid-training TD(0) evaluation on a $7\times 7$ board (10k self-play games, hidden size 256). \textbf{Top:} directional and aggregated win rates versus random and heuristic agents. \textbf{Bottom:} checkpoint outcomes for current versus previous (Lost both, Won only as P2, Won only as P1, Won both). The dotted line marks checkpoint 1900 (selected best checkpoint), using a rough composite score $\mathrm{Score} = w_{\mathrm{prev}}\,\mathrm{WR}(\mathrm{curr\ vs\ prev}) + w_{\mathrm{heur}}\,\mathrm{WR}(\mathrm{curr\ vs\ heuristic}) + w_{\mathrm{rand}}\,\mathrm{WR}(\mathrm{curr\ vs\ random})$. Across runs, changing hidden size, number of training games, or using the best checkpoint did not show a consistent impact on performance.}
  \label{fig:training_output}
\end{figure}

\begin{figure}[h]
  \centering
  \includegraphics[width=\columnwidth]{sb3_elo_vs_iters.png}
  \caption{SB3 agent Elo rating vs.\ cumulative training iterations (500\,k
    steps per stage, 4 stages). Each stage boundary is marked with a vertical
    dashed line. Shaded region indicates $\pm1$ standard error. MaskablePPO
    improves monotonically; PPO degrades after stage 2, consistent with
    self-play induced policy collapse.}
  \label{fig:sb3_elo}
\end{figure}

\begin{figure}[h]
  \centering
  \includegraphics[width=\columnwidth]{heatmap.png}
  \caption{Head-to-head win-rate heatmap (combined across both seat assignments,
    $N=100$ games per unordered pair, main tournament). Darker cells indicate
    higher win rate for the row agent. The block-diagonal structure confirms three
    distinct performance tiers.}
  \label{fig:heatmap}
\end{figure}

% ══════════════════════════════════════════════════════════════════════════
\section{Discussion}
% ══════════════════════════════════════════════════════════════════════════

\subsection{Search Is the Dominant Factor}

The ${\sim}150$ Elo gap between the best non-search agent (MaskablePPO, 1740)
and the weakest search agent (MCTS(500), 1869) is a significant finding
in this study. Runtime tree search provides a form of implicit planning that
learned value functions cannot replicate at comparable compute budgets. This is
consistent with the broader literature showing that search scales more efficiently
than network depth for combinatorial games at moderate board sizes \cite{silver2017}.

\subsection{Action Masking Is Critical for Deep RL}

The 371 Elo gap between MaskablePPO and PPO despite identical network
architecture, training budget, and self-play curriculum demonstrates that action
masking is the single most impactful design decision for applying deep RL to board
games. On a 28-cell board with an average of 18 legal moves, roughly 36\% of the
policy gradient in standard PPO is wasted on illegal cells. Masking eliminates
this waste entirely and is essentially free to implement using sb3-contrib.

\subsection{Self-Play Curriculum Effects Diverge by Algorithm}

Self-play consistently benefits search-based agents (PV-MCTS, SP-PV-MCTS)
and mildly helps stable on-policy agents (A2C), but harms PPO. The instability
arises because PPO's clipping mechanism was designed for stationary environments;
when the opponent distribution shifts every 500\,k steps, the old-policy
reference used by the clip objective becomes stale faster than the clipping
horizon compensates. This motivates opponent-pool methods such as PFSP
\cite{vinyals2019} or league training for future work.

\subsection{TD Agents Are Fundamentally Limited by Lookahead Depth}

The near-identical performance of TD(0) and TD($\lambda$) indicates that credit
assignment is not the bottleneck for TD agents in Y. The bottleneck is lookahead
depth: 1-ply evaluation cannot capture the multi-move threats that characterise
Y strategy. Extending TD agents to 2-ply or 3-ply search would be a direct test
of this hypothesis.

% ══════════════════════════════════════════════════════════════════════════
\section{Conclusion}
% ══════════════════════════════════════════════════════════════════════════

This was our comprehensive RL study of the connection game Y,
implementing and evaluating nine agents across three paradigms on a size-7 board.
Our main findings are:

\begin{enumerate}
  \item \textbf{Search dominates:} MCTS and its policy/value-guided variants
    occupy the top four positions with a ${\sim}500$ Elo gap to all non-search
    methods. Runtime lookahead cannot be replaced by learned value functions at
    modest compute budgets.
  \item \textbf{Action masking is decisive for deep RL:} MaskablePPO outperforms
    all TD agents and matches the hand-coded Heuristic after 2\,M steps.
    Plain PPO without masking performs near the TD tier.
  \item \textbf{Algorithm stability matters for self-play:} A2C outperforms PPO
    under self-play due to its simpler, more conservative update rule. Self-play
    curriculum design is critical for on-policy gradient methods.
  \item \textbf{TD methods share a representational ceiling:} TD(0), TD($\lambda$),
    and TD-CNN cluster within 160 Elo points of each other and of the no-search
    CNN baseline, confirming that value function quality is not the binding
    constraint.
\end{enumerate}

\subsection{Future Work}

Several extensions are immediately motivated by these results:
\begin{itemize}
  \item \textbf{Larger boards} (size 9–19): The performance hierarchy may
    change as the branching factor grows and rollout quality becomes less reliable
    for pure MCTS.
  \item \textbf{MaskablePPO + MCTS guidance:} Training MaskablePPO from MCTS
    visit-count targets (as in AlphaZero) rather than environment rewards could
    close the gap between the SB3 and search tiers.
  \item \textbf{Opponent-pool self-play:} PFSP or league training for PPO-based
    agents to prevent policy collapse under non-stationary opponents.
  \item \textbf{Multi-ply TD lookahead:} Extending 1-ply TD to 2- or 3-ply to
    isolate whether search depth or value accuracy is the binding constraint.
  \item \textbf{Continuous-action agents:} DDPG, TD3, and SAC with the masked
    argmax environment wrapper remain untested and could provide interesting
    comparisons against discrete-action baselines.
\end{itemize}



% ══════════════════════════════════════════════════════════════════════════
% References
% ══════════════════════════════════════════════════════════════════════════
\begin{thebibliography}{99}

\bibitem{silver2016}
D. Silver et al., ``Mastering the game of Go with deep neural networks and tree
search,'' \textit{Nature}, vol.\ 529, pp.\ 484--489, 2016.

\bibitem{silver2017}
D. Silver et al., ``Mastering Chess and Shogi by Self-Play with a General
Reinforcement Learning Algorithm,'' \textit{arXiv:1712.01815}, 2017.

\bibitem{schrittwieser2021}
J. Schrittwieser et al., ``Online and Offline Reinforcement Learning by Planning
with a Learned Model,'' \textit{arXiv:2104.06294}, 2021.

\bibitem{anthony2017}
T. Anthony, Z. Tian, and D. Barber, ``Thinking Fast and Slow with Deep Learning
and Tree Search,'' \textit{Advances in Neural Information Processing Systems},
vol.\ 30, 2017.

\bibitem{harbecke2025}
D. Harbecke, ``HexHex: AlphaGo Zero adaptation for Hex,'' GitHub repository,
2025. [Online]. Available: \url{https://github.com/harbecke/HexHex}

\bibitem{tesauro1995}
G. Tesauro, ``Temporal difference learning and TD-Gammon,''
\textit{Communications of the ACM}, vol.\ 38, no.\ 3, pp.\ 58--68, 1995.

\bibitem{campbell2002}
M. Campbell, A. J. Hoane, and F.-h. Hsu, ``Deep Blue,'' \textit{Artificial
Intelligence}, vol.\ 134, pp.\ 57--83, 2002.

\bibitem{mnih2015}
V. Mnih et al., ``Human-level control through deep reinforcement learning,''
\textit{Nature}, vol.\ 518, pp.\ 529--533, 2015.

\bibitem{schulman2017}
J. Schulman, F. Wolski, P. Dhariwal, A. Radford, and O. Klimov, ``Proximal
Policy Optimization Algorithms,'' \textit{arXiv:1707.06347}, 2017.

\bibitem{raffin2021}
A. Raffin et al., ``Stable-Baselines3: Reliable Reinforcement Learning
Implementations,'' \textit{Journal of Machine Learning Research}, vol.\ 22,
no.\ 268, pp.\ 1--8, 2021.

\bibitem{sb3contrib}
A. Raffin et al., ``SB3-Contrib: Contrib Package for Stable-Baselines3,''
GitHub repository, 2021. [Online]. Available:
\url{https://github.com/Stable-Baselines-Team/stable-baselines3-contrib}

\bibitem{bradleyterry1952}
R. A. Bradley and M. E. Terry, ``Rank analysis of incomplete block designs:
I. The method of paired comparisons,'' \textit{Biometrika}, vol.\ 39,
pp.\ 324--345, 1952.

\bibitem{veness2009}
J. Veness, D. Silver, W. Uther, and A. Blair, ``Bootstrapping from Game Tree
Search,'' \textit{Advances in Neural Information Processing Systems}, 2009.

\bibitem{balduzzi2019}
D. Balduzzi et al., ``Open-ended learning leads to generally capable agents,''
\textit{arXiv:1901.08568}, 2019.

\bibitem{vinyals2019}
O. Vinyals et al., ``Grandmaster level in StarCraft II using multi-agent
reinforcement learning,'' \textit{Nature}, vol.\ 575, pp.\ 350--354, 2019.

\end{thebibliography}

\end{document}